\documentclass[12pt, a4paper]{article}
\usepackage{../../../myconfig} % Gọi file cấu hình từ ngoài gốc vào

% ==========================================
% BẮT ĐẦU NỘI DUNG TÀI LIỆU
% ==========================================
\begin{document}

% Truyền 3 tham số: {Nhãn cam} {Chữ to chính giữa} {Chữ ở góc phải Header}
\titlebox{Chủ đề: Tích phân}{Tích phân}{Nguyên hàm - Tích phân: Tích phân}

% --- CÂU 1 ---
\questionLinesManual{6}{1}{Biết \( f(x) \) là hàm số liên tục trên tập \( \mathbb{R} \), \( a \) là một số thực thỏa mãn \( 0 < a < \pi \),     \( \displaystyle \int_{0}^{a} f(x) \, dx = \int_{\pi}^{a} f(x) \, dx = 1 \). Tính \( \displaystyle \int_{0}^{\pi} f(x) \, dx \).}{
    \begin{tasks}(4)
        \task $0$
        \task $2$
        \task $\dfrac{1}{2}$
        \task $1$
    \end{tasks}
}

% --- CÂU 2 ---
\questionLinesManual{6}{2}{Cho \( \displaystyle \int_{0}^{1} f(x) \, dx = 1 \) và \( \displaystyle \int_{0}^{2} f(x) \, dx = -4 \). Tích phân \( \displaystyle \int_{1}^{2} f(x) \, dx \) bằng:}{
    \begin{tasks}(4)
        \task $5$
        \task $-3$
        \task $-5$
        \task $3$
    \end{tasks}
}

% --- CÂU 3 ---
\questionLinesManual{6}{3}{Cho \( I = \displaystyle \int_{0}^{2} f(x) \, dx = 3 \). Khi đó \( J = \displaystyle \int_{0}^{2} [4f(x) - 3] \, dx \) bằng:}{
    \begin{tasks}(4)
        \task $8$
        \task $6$
        \task $2$
        \task $4$
    \end{tasks}
}

% --- CÂU 4 ---
\questionLinesManual{6}{4}{Cho hàm số \( f(x) \) liên tục trên \([a; c]\) và \( b \) là số thực tùy ý thuộc đoạn \([a; c]\). Nếu biết \( \displaystyle \int_a^b f(x) \, dx = -5 \) và \( \displaystyle \int_b^c f(x) \, dx = 10 \), thì giá trị của \( \displaystyle \int_a^c f(x) \, dx \) là bao nhiêu?}{
    \begin{tasks}(4)
        \task $5$
        \task $-5$
        \task $15$
        \task $-15$
    \end{tasks}
}

% --- CÂU 5 ---
\questionLinesManual{6}{5}{Một vật chuyển động với vận tốc \( v(t) = 4t + 8 \, (m/s) \), với thời gian \( t \) tính bằng giây. Tính quãng đường vật đi được trong khoảng thời gian từ \( t = 8 \) đến \( t = 10 \).}{
    \begin{tasks}(4)
        \task $89 \, (m)$
        \task $87 \, (m)$
        \task $86 \, (m)$
        \task $88 \, (m)$
    \end{tasks}
}

% --- CÂU 6 ---
\questionLinesManual{6}{6}{Cho hàm số \( y = f(x) \) có đạo hàm liên tục trên đoạn \([a; b]\) và \( f(a) = -1, \, f(b) = 3 \). Khi đó \( \displaystyle \int_a^b f'(x) \, dx \) bằng:}{
    \begin{tasks}(4)
        \task $-3$
        \task $4$
        \task $-4$
        \task $2$
    \end{tasks}
}

% --- CÂU 7 ---
\questionLinesManual{6}{7}{Biết \( \displaystyle \int_{0}^{9} f(x) \, dx = 37 \) và \( \displaystyle \int_{0}^{9} [3f(x) - 2g(x)] \, dx = 61 \). Tính \( \displaystyle \int_{0}^{9} g(x) \, dx \) bằng}{
    \begin{tasks}(4)
        \task $-25$
        \task $25$
        \task $-86$
        \task $86$
    \end{tasks}
}

% --- CÂU 8 ---
\questionLinesManual{6}{8}{Cho hàm số \( y = f(x) \) liên tục trên đoạn \( [a; b] \) và \( c \) là số thực tùy ý thuộc đoạn \( [a; b] \). Nếu \( \displaystyle \int_{a}^{b} f(x) \, dx = 3 \) và \( \displaystyle \int_{a}^{c} f(x) \, dx = 8 \) thì tích phân \( \displaystyle \int_{c}^{b} f(x) \, dx \) bằng}{
    \begin{tasks}(4)
        \task $11$
        \task $-5$
        \task $5$
        \task $-11$
    \end{tasks}
}

% --- CÂU 9 ---
\questionLinesManual{6}{9}{Biết rằng \( F(x) \) là một nguyên hàm của hàm số \( f(x) \) trên đoạn \( [1; 4] \) và \( F(4) = 9, \ F(1) = 3 \). Giá trị của \( \displaystyle \int_{1}^{4} [f(x) + 2] \, dx \) bằng}{
    \begin{tasks}(4)
        \task $0$
        \task $8$
        \task $-4$
        \task $12$
    \end{tasks}
}

% --- CÂU 10 ---
\questionLinesManual{6}{10}{Nếu \( \displaystyle \int_{-3}^{1} f(x) \, dx = -2 \) thì \( \displaystyle \int_{-3}^{1} [2 - 5f(x)] \, dx \) bằng}{
    \begin{tasks}(4)
        \task $18$
        \task $6$
        \task $12$
        \task $-4$
    \end{tasks}
}

% --- CÂU 11 ---
\questionLinesManual{6}{11}{Cho \( \displaystyle \int_{0}^{2} [f(x) - 3x^2] \, dx = 4 \). Tích phân \( \displaystyle \int_{0}^{2} f(x) \, dx \) bằng}{
    \begin{tasks}(4)
        \task $8$
        \task $-4$
        \task $12$
        \task $4$
    \end{tasks}
}

% --- CÂU 12 ---
\questionLinesManual{6}{12}{Cho hàm số \( y = f(x) \) liên tục trên \( \mathbb{R} \). Biết \( \displaystyle \int_{5}^{9} f(x) \, dx = 25 \) thì \( \displaystyle \int_{5}^{9} f(x) \, dx \) bằng}{
    \begin{tasks}(4)
        \task $9$
        \task $25$
        \task $-25$
        \task $5$
    \end{tasks}
}

% --- CÂU 13 ---
\questionLinesManual{6}{13}{Nếu \( \displaystyle \int_{0}^{4} f(x) \, dx = 6 \) và \( \displaystyle \int_{1}^{4} f(x) \, dx = -5 \) thì \( \displaystyle \int_{0}^{1} f(x) \, dx \) bằng:}{
    \begin{tasks}(4)
        \task $11$
        \task $-11$
        \task $-1$
        \task $1$
    \end{tasks}
}

% --- CÂU 14 ---
\questionLinesManual{6}{14}{Nếu \( \displaystyle \int_{0}^{2} f(x) \, dx = 3 \) thì \( \displaystyle \int_{0}^{2} [f(x) + 2] \, dx \) bằng:}{
    \begin{tasks}(4)
        \task $5$
        \task $6$
        \task $7$
        \task $10$
    \end{tasks}
}

% --- CÂU 15 ---
\questionLinesManual{6}{15}{Cho hàm số \( y = f(x) \) liên tục trên \( \mathbb{R} \). Biết hàm số \( F(x) \) là một nguyên hàm của hàm số \( f(x) \) trên \( \mathbb{R} \) thoả mãn \( F(5) = 2 + F(1) \). Giá trị của \( \displaystyle \int_{1}^{5} f(x) \, dx \) là:}{
    \begin{tasks}(4)
        \task $8$
        \task $2$
        \task $-2$
        \task $-8$
    \end{tasks}
}

% --- CÂU 16 ---
\questionLinesManual{6}{16}{Biết \( \displaystyle \int_{0}^{1} [f(x) + 2x] \, dx = 5 \). Khi đó \( \displaystyle \int_{0}^{1} f(x) \, dx \) bằng}{
    \begin{tasks}(4)
        \task $7$
        \task $3$
        \task $5$
        \task $4$
    \end{tasks}
}

% --- CÂU 17 ---
\questionLinesManual{6}{17}{Cho \( \displaystyle \int_{1}^{2} f(x) \, dx = -3 \) và \( \displaystyle \int_{1}^{2} g(x) \, dx = 4 \). Giá trị tích phân \( \displaystyle \int_{1}^{2} (g(x) + 2f(x)) \, dx \) bằng}{
    \begin{tasks}(4)
        \task $-2$
        \task $2$
        \task $1$
        \task $5$
    \end{tasks}
}

% --- CÂU 18 ---
\questionLinesManual{6}{18}{Nếu \( \displaystyle \int_{0}^{3} f(x) \, dx = 6 \) thì \( \displaystyle \int_{0}^{3} \left[ \dfrac{1}{3} f(x) + 2 \right] \, dx \) bằng:}{
    \begin{tasks}(4)
        \task $8$
        \task $6$
        \task $5$
        \task $9$
    \end{tasks}
}

% --- CÂU 19 ---
\questionLinesManual{6}{19}{Cho hàm số \( f(x) \) liên tục trên \( \mathbb{R} \) thỏa mãn \( \displaystyle \int_{0}^{1} f(x) \, dx = 2 \) và \( \displaystyle \int_{0}^{2} f(x) \, dx = 7 \). Khi đó \( \displaystyle \int_{1}^{2} f(x) \, dx \) bằng}{
    \begin{tasks}(4)
        \task $-9$
        \task $9$
        \task $-5$
        \task $5$
    \end{tasks}
}

% --- CÂU 20 ---
\questionLinesManual{6}{20}{Nếu \( \displaystyle \int_{0}^{\pi/2} f(x) \, dx = -5 \) thì \( \displaystyle \int_{0}^{\pi/2} [f(x) + \sin x] \, dx \) bằng}{
    \begin{tasks}(4)
        \task $-4$
        \task $-7$
        \task $-6$
        \task $-3$
    \end{tasks}
}

% --- CÂU 21 ---
\questionLinesManual{6}{21}{Cho hàm số \( f(x) \) liên tục trên \( \mathbb{R} \) và có \( \displaystyle \int_{0}^{4} f(x) \, dx = 9 \), \( \displaystyle \int_{0}^{4} g(x) \, dx = 4 \). Tích phân \( \displaystyle \int_{0}^{4} [2f(x) + 3g(x)] \, dx \) bằng}{
    \begin{tasks}(4)
        \task $13$
        \task $5$
        \task $\dfrac{9}{4}$
        \task $36$
    \end{tasks}
}

% --- CÂU 22 ---
\questionLinesManual{6}{22}{Nếu \( \displaystyle \int_{0}^{2} f(x) \, dx = 4 \) thì \( \displaystyle \int_{0}^{1} \left[ \dfrac{1}{2} f(x) - 2 \right] dx \) bằng}{
    \begin{tasks}(4)
        \task $4$
        \task $6$
        \task $0$
        \task $-2$
    \end{tasks}
}

% --- CÂU 23 ---
\questionLinesManual{6}{23}{Biết \( F(x) = x^3 \) là một nguyên hàm của hàm số \( f(x) \) trên \( \mathbb{R} \). Giá trị của \( \displaystyle \int_{1}^{2} [2 + f(x)] \, dx \) bằng}{
    \begin{tasks}(4)
        \task $9$
        \task $\dfrac{15}{4}$
        \task $7$
        \task $\dfrac{23}{4}$
    \end{tasks}
}

% --- CÂU 24 ---
\questionLinesManual{6}{24}{Nếu \( \displaystyle \int_{a}^{b} f(x) \, dx = 2025 \) thì \( \displaystyle \int_{a}^{b} 2f(x) \, dx \) bằng?}{
    \begin{tasks}(4)
        \task $2025^2$
        \task $\dfrac{2025}{2}$
        \task $2023$
        \task $4050$
    \end{tasks}
}

% --- CÂU 25 ---
\questionLinesManual{6}{25}{Cho \( \displaystyle \int_{0}^{1} f(x) \, dx = 7 \). Giá trị của \( \displaystyle \int_{0}^{1} 3f(x) \, dx \) bằng:}{
    \begin{tasks}(4)
        \task $21$
        \task $\dfrac{3}{7}$
        \task $7$
        \task $\dfrac{7}{3}$
    \end{tasks}
}

% --- CÂU 26 ---
\questionLinesManual{6}{26}{Cho hàm số \( F(x) \) là nguyên hàm của \( f(x) \). Biết \( F(1) = -3, F(-2) = 12 \). Tính \( I = \displaystyle \int_{-2}^{1} f(x) \, dx \).}{
    \begin{tasks}(4)
        \task $15$
        \task $-15$
        \task $9$
        \task $-36$
    \end{tasks}
}

% --- CÂU 27 ---
\questionLinesManual{6}{27}{Cho hàm số \( f(x) \) liên tục trên \([a; b]\) và có một nguyên hàm trên \([a; b]\) là hàm số \( F(x) \). Tìm mệnh đề đúng trong các mệnh đề sau}{
    \begin{tasks}(2)
        \task \( \displaystyle \int_a^b f(x) \, dx = f(b) - f(a) \)
        \task \( \displaystyle \int_a^b f(x) \, dx = F(a) - F(b) \)
        \task \( \displaystyle \int_a^b f(x) \, dx = F(b) - F(a) \)
        \task \( \displaystyle \int_a^b f(x) \, dx = f(a) - f(b) \)
    \end{tasks}
}

% --- CÂU 28 ---
\questionLinesManual{6}{28}{Cho hàm số \( f(x) \) có đạo hàm \( f'(x) \) liên tục trên \( \mathbb{R} \), \( f(1) = 1 \) và \( \displaystyle \int_1^4 f'(x) \, dx = 15 \) khi đó giá trị \( f(4) \) bằng:}{
    \begin{tasks}(4)
        \task $16$
        \task $15$
        \task $14$
        \task $17$
    \end{tasks}
}

% --- CÂU 29 ---
\questionLinesManual{6}{29}{Phát biểu nào sau đây là đúng?}{
    \begin{tasks}(2)
        \task \( \displaystyle \int_a^b x^\alpha \, dx = b^{\alpha+1} - a^{\alpha+1} \)
        \task \( \displaystyle \int_a^b x^\alpha \, dx = \alpha \left( b^{\alpha-1} - a^{\alpha-1} \right) \)
        \task \( \displaystyle \int_a^b x^\alpha \, dx = \dfrac{b^{\alpha+1} - a^{\alpha+1}}{\alpha + 1} \quad (\alpha \neq -1) \)
        \task \( \displaystyle \int_a^b x^\alpha \, dx = \dfrac{b^{\alpha+1} - a^{\alpha+1}}{\alpha} \quad (\alpha \neq 0) \)
    \end{tasks}
}

% --- CÂU 30 ---
\questionLinesManual{6}{30}{Phát biểu nào sau đây là đúng?}{
    \begin{tasks}(2)
        \task \( \displaystyle \int_a^b \sin x \, dx = \sin a - \sin b \)
        \task \( \displaystyle \int_a^b \sin x \, dx = \sin b - \sin a \)
        \task \( \displaystyle \int_a^b \sin x \, dx = \cos a - \cos b \)
        \task \( \displaystyle \int_a^b \sin x \, dx = \cos b - \cos a \)
    \end{tasks}
}

% --- CÂU 31 ---
\questionLinesManual{6}{31}{Phát biểu nào sau đây là đúng? Biết \( f(x) = \dfrac{1}{\sin^2 x} \) liên tục trên \([a; b]\).}{
    \begin{tasks}(2)
        \task \( \displaystyle \int_a^b \dfrac{1}{\sin^2 x} \, dx = \cot a - \cot b \)
        \task \( \displaystyle \int_a^b \dfrac{1}{\sin^2 x} \, dx = -\cot b - \cot a \)
        \task \( \displaystyle \int_a^b \dfrac{1}{\sin^2 x} \, dx = \tan a - \tan b \)
        \task \( \displaystyle \int_a^b \dfrac{1}{\sin^2 x} \, dx = -\tan b - \tan a \)
    \end{tasks}
}

% --- CÂU 32 ---
\questionLinesManual{6}{32}{Phát biểu nào sau đây là đúng?}{
    \begin{tasks}(2)
        \task \( \displaystyle \int_a^b e^x \, dx = e^{b+1} - e^{a+1} \)
        \task \( \displaystyle \int_a^b e^x \, dx = e^{a+1} - e^{b+1} \)
        \task \( \displaystyle \int_a^b e^x \, dx = e^{b} - e^{a} \)
        \task \( \displaystyle \int_a^b e^x \, dx = e^{a} - e^{b} \)
    \end{tasks}
}

% --- CÂU 33 ---
\questionLinesManual{6}{33}{Tích phân \( \displaystyle \int_a^b \dfrac{1}{x} \, dx \) bằng:}{
    \begin{tasks}(4)
        \task \( \ln b - \ln a \)
        \task \( |\ln b| - |\ln a| \)
        \task \( \ln |b| - \ln |a| \)
        \task \( \ln |a| - \ln |b| \)
    \end{tasks}
}

% --- CÂU 34 ---
\questionLinesManual{15}{34}{Tích phân \( \displaystyle \int_{0}^{4} |x^2 - 4| \, dx \) bằng}{
    \begin{tasks}(4)
        \task $16$
        \task $9$
        \task $5$
        \task $6$
    \end{tasks}
}

% --- CÂU 35 ---
\questionLinesManual{10}{35}{Tích phân \( \displaystyle \int_{-2}^{1} |2x + 2| \, dx \) bằng}{
    \begin{tasks}(4)
        \task $12$
        \task $9$
        \task $5$
        \task $6$
    \end{tasks}
}

% --- CÂU 36 ---
\questionLinesManual{10}{36}{Cho \( F(x) \) là một nguyên hàm của \( f(x) = \dfrac{2}{x+2} \). Biết \( F(-1) = 0 \). Tính \( F(2) \) kết quả là.}{
    \begin{tasks}(4)
        \task \( \ln 8 + 1 \)
        \task \( 4\ln 2 + 1 \)
        \task \( 2\ln 3 + 2 \)
        \task \( 2\ln 4 \)
    \end{tasks}
}

% --- CÂU 37 ---
\questionLinesManual{10}{37}{Cho hàm số \( f(x) \). Biết \( f(0) = 4 \) và \( f'(x) = 2\sin^2 x + 1, \, \forall x \in \mathbb{R} \), khi đó \( \displaystyle \int_{0}^{\frac{\pi}{4}} f(x) \, dx \) bằng}{
    \begin{tasks}(4)
        \task \( \dfrac{\pi^2 + 16\pi - 4}{16} \)
        \task \( \dfrac{\pi^2 - 4}{16} \)
        \task \( \dfrac{\pi^2 + 15\pi}{16} \)
        \task \( \dfrac{\pi^2 + 16\pi - 16}{16} \)
    \end{tasks}
}

% --- CÂU 38 ---
\questionLinesManual{10}{38}{Cho \( \displaystyle \int_{3}^{4} \dfrac{5x - 8}{x^2 - 3x + 2} \, dx = a\ln 3 + b\ln 2 + c\ln 5 \), với \( a, b, c \) là các số hữu tỉ. Giá trị của \( 2^{a - 3b + c} \) bằng}{
    \begin{tasks}(4)
        \task $12$
        \task $6$
        \task $1$
        \task $64$
    \end{tasks}
}

% --- CÂU 39 ---
\questionLinesManual{10}{39}{Biết \( \displaystyle \int_{3}^{5} \dfrac{x^2 + x + 1}{x + 1} \, dx = a + \ln \dfrac{b}{2} \) với \( a, b \) là các số nguyên. Tính \( S = a - 2b \).}{
    \begin{tasks}(4)
        \task $S = 2$
        \task $S = -2$
        \task $S = 5$
        \task $S = 10$
    \end{tasks}
}

% --- CÂU 40 ---
\questionLinesManual{10}{40}{Biết rằng \( \displaystyle \int_{0}^{1} \dfrac{1}{x^2 + x + 1} \, dx = \dfrac{\pi\sqrt{a}}{b} \quad (a, b \in \mathbb{Z}, a < 10) \). Khi đó \( a + b \) có giá trị bằng}{
    \begin{tasks}(4)
        \task $14$
        \task $15$
        \task $13$
        \task $12$
    \end{tasks}
}

% --- CÂU 41 ---
\questionLinesManual{10}{41}{Biết \( \displaystyle \int_{0}^{1} \dfrac{x^2 + 5x + 2}{x^2 + 4x + 3} \, dx = a + b\ln 3 + c\ln 5 \), \( (a, b, c \in \mathbb{Q}) \). Giá trị của \( abc \) bằng}{
    \begin{tasks}(4)
        \task $-8$
        \task $-10$
        \task $-12$
        \task $16$
    \end{tasks}
}

% --- CÂU 42 ---
\questionLinesManual{10}{42}{Cho biết \( \displaystyle \int_{0}^{2} \dfrac{x-1}{x^2+4x+3} \, dx = a\ln 5 + b\ln 3 \), với \( a, b \in \mathbb{Q} \). Tính \( T = a^2 + b^2 \) bằng}{
    \begin{tasks}(4)
        \task $13$
        \task $10$
        \task $25$
        \task $5$
    \end{tasks}
}

% --- CÂU 43 ---
\questionLinesManual{10}{43}{Cho \( \displaystyle \int_{\frac{1}{2}}^{3} \dfrac{dx}{(x+1)(x+2)} = a\ln 2 + b\ln 3 + c\ln 5 \) với \( a, b, c \) là các số hữu tỉ. Giá trị của \( a + b^2 - c^3 \) bằng}{
    \begin{tasks}(4)
        \task $3$
        \task $6$
        \task $5$
        \task $4$
    \end{tasks}
}

% --- CÂU 44 ---
\questionLinesManual{10}{44}{Có bao nhiêu giá trị nguyên dương của \( a \) để \( \displaystyle \int_{0}^{a} (2x - 3) \, dx \leq 4 \)?}{
    \begin{tasks}(4)
        \task $5$
        \task $6$
        \task $4$
        \task $3$
    \end{tasks}
}

% --- CÂU 45 ---
\questionLinesManual{10}{45}{Cho \( \displaystyle \int_{0}^{m} (3x^2 - 2x + 1) \, dx = 6 \). Giá trị của tham số \( m \) thuộc khoảng nào sau đây?}{
    \begin{tasks}(4)
        \task $(-1; 2)$
        \task $(-\infty; 0)$
        \task $(0; 4)$
        \task $(-3; 1)$
    \end{tasks}
}

% --- CÂU 46: TÍNH CHẤT TÍCH PHÂN ---
\questionTrueFalse{8}{46}{Giả sử \( f \) là hàm số liên tục trên khoảng \( K \) và \( a, b, c \) là ba số bất kỳ trên khoảng \( K \). Các mệnh đề sau đúng hay sai?}{
    \textbf{a)} & \( \displaystyle \int_a^b f(x) \, dx = 1 \).
    & \textcolor{amberorange}{$\bigcirc$} & \textcolor{amberorange}{$\bigcirc$} \\ \hline
    
    \textbf{b)} & \( \displaystyle \int_a^b f(x) \, dx = -\int_b^a f(x) \, dx \).
    & \textcolor{amberorange}{$\bigcirc$} & \textcolor{amberorange}{$\bigcirc$} \\ \hline
    
    \textbf{c)} & \( \displaystyle \int_a^c f(x) \, dx + \int_c^b f(x) \, dx = \int_a^b f(x) \, dx, \, c \in (a; b) \).
    & \textcolor{amberorange}{$\bigcirc$} & \textcolor{amberorange}{$\bigcirc$} \\ \hline
    
    \textbf{d)} & \( \displaystyle \int_a^b x f(x) \, dx = x \int_a^b f(x) \, dx \).
    & \textcolor{amberorange}{$\bigcirc$} & \textcolor{amberorange}{$\bigcirc$} \\
}

% --- CÂU 47: TÍCH PHÂN HÀM BẬC 2 ---
\questionTrueFalse{15}{47}{Cho hàm số \( f(x) = ax^2 + bx + 1 \) (\( a, b \in \mathbb{R} \)) thoả mãn \( \displaystyle \int_0^1 f(x) \, dx = 4 \) và \( \displaystyle \int_1^2 f(x) \, dx = 14 \). Các mệnh đề sau đây đúng hay sai?}{
    \textbf{a)} & \( \displaystyle \int_0^1 f(x) \, dx = -14 \).
    & \textcolor{amberorange}{$\bigcirc$} & \textcolor{amberorange}{$\bigcirc$} \\ \hline
    
    \textbf{b)} & \( \displaystyle \int_0^2 f(x) \, dx = 18 \).
    & \textcolor{amberorange}{$\bigcirc$} & \textcolor{amberorange}{$\bigcirc$} \\ \hline
    
    \textbf{c)} & \( a + b = 6 \).
    & \textcolor{amberorange}{$\bigcirc$} & \textcolor{amberorange}{$\bigcirc$} \\ \hline
    
    \textbf{d)} & \( \displaystyle \int_0^3 f(x) \, dx = 44 \).
    & \textcolor{amberorange}{$\bigcirc$} & \textcolor{amberorange}{$\bigcirc$} \\
}

% --- CÂU 48: NGUYÊN HÀM HÀM PHÂN THỨC ---
\questionTrueFalse{15}{48}{Cho hàm số \( F(x) \) là một nguyên hàm của hàm số \( f(x) = \dfrac{2x^3 - 2x - 1}{x^2} \), với mọi \( x \neq 0 \).}{
    \textbf{a)} & \( F'(x) = \dfrac{2x^3 - 2x - 1}{x^2} \).
    & \textcolor{amberorange}{$\bigcirc$} & \textcolor{amberorange}{$\bigcirc$} \\ \hline
    
    \textbf{b)} & \( \displaystyle \int f(x) \, dx = x^2 - 2\ln|x| - \dfrac{1}{x} + C \) với \( C \) là hằng số.
    & \textcolor{amberorange}{$\bigcirc$} & \textcolor{amberorange}{$\bigcirc$} \\ \hline
    
    \textbf{c)} & Nếu \( F(1) = -1 \) thì \( F(3) = \dfrac{19}{3} - 2\ln 3 \).
    & \textcolor{amberorange}{$\bigcirc$} & \textcolor{amberorange}{$\bigcirc$} \\ \hline
    
    \textbf{d)} & Biết rằng \( \displaystyle \int_1^2 f(x) \, dx = \dfrac{a}{b} - \ln c \) với \( a, b, c \in \mathbb{N}^*, \dfrac{a}{b} \) là phân số tối giản. Khi đó \( a + b - c = 3 \).
    & \textcolor{amberorange}{$\bigcirc$} & \textcolor{amberorange}{$\bigcirc$} \\
}

% --- CÂU 49: VẬN TỐC VÀ GIA TỐC ---
\questionTrueFalse{15}{49}{Giả sử \( v(t) \) là phương trình vận tốc của một vật chuyển động theo thời gian \( t \) (giây), \( a(t) \) là phương trình gia tốc của vật đó chuyển động theo thời gian \( t \) (giây). Xét chuyển động trong khoảng thời gian từ \( c \) (giây) đến \( b \) (giây). Các mệnh đề sau đúng hay sai?}{
    \textbf{a)} & \( \displaystyle \int_c^b a(t) \, dt = v(b) - v(c) \).
    & \textcolor{amberorange}{$\bigcirc$} & \textcolor{amberorange}{$\bigcirc$} \\ \hline
    
    \textbf{b)} & \( \displaystyle \int_c^b v(t) \, dt = a(b) - a(c) \).
    & \textcolor{amberorange}{$\bigcirc$} & \textcolor{amberorange}{$\bigcirc$} \\ \hline
    
    \textbf{c)} & \( \displaystyle \int_c^b v'(t) \, dt = v(c) - v(b) \).
    & \textcolor{amberorange}{$\bigcirc$} & \textcolor{amberorange}{$\bigcirc$} \\ \hline
    
    \textbf{d)} & \( \displaystyle \int_c^b v'(t) \, dt = v(b) - v(c) \).
    & \textcolor{amberorange}{$\bigcirc$} & \textcolor{amberorange}{$\bigcirc$} \\
}

% --- CÂU 50: ỨNG DỤNG TÍCH PHÂN ---
\questionTrueFalse{15}{50}{Một người điều khiển ô tô đang ở đường dẫn muốn nhập làn vào đường cao tốc. Khi xe ô tô cách điểm nhập làn 200m, tốc độ của ô tô là \( 36 \, \text{km/h} \). Hai giây sau đó, ô tô bắt đầu tăng tốc với tốc độ \( v(t) = at + b \) (\( a, b \in \mathbb{R}, a > 0 \)), trong đó \( t \) là thời gian tính bằng giây kể từ khi bắt đầu tăng tốc. Biết rằng ô tô nhập làn cao tốc sau 12 giây và duy trì sự tăng tốc đó trong 24 giây kể từ khi bắt đầu tăng tốc.}{
    \textbf{a)} & Quãng đường ô tô đi được từ khi bắt đầu tăng tốc đến khi nhập làn là \( 180 \, \text{m} \).
    & \textcolor{amberorange}{$\bigcirc$} & \textcolor{amberorange}{$\bigcirc$} \\ \hline
    
    \textbf{b)} & Giá trị của \( b = 10 \).
    & \textcolor{amberorange}{$\bigcirc$} & \textcolor{amberorange}{$\bigcirc$} \\ \hline
    
    \textbf{c)} & Quãng đường \( S(t) \) (đơn vị: mét) mà ô tô đi được trong thời gian \( t \) giây (\( 0 \leq t \leq 24 \)) kể từ khi tăng tốc được tính theo công thức  
    \[S(t) = \int_{0}^{24} v(t) \, dt.\]
    & \textcolor{amberorange}{$\bigcirc$} & \textcolor{amberorange}{$\bigcirc$} \\ \hline
    
    \textbf{d)} & Sau 24 giây kể từ khi tăng tốc, tốc độ của ô tô không vượt quá tốc độ tối đa cho phép là \( 100 \, \text{km/h} \).
    & \textcolor{amberorange}{$\bigcirc$} & \textcolor{amberorange}{$\bigcirc$} \\
}
\end{document}
